\chapter{Funkcje pierwotne}
\label{ma1:ch:primitives}

Dotychczas wyznaczaliśmy pochodną danej funkcji. Teraz odwracamy to zadanie: dla zadanej funkcji $f$ szukamy funkcji $F$, której pochodną jest $f$. Rozróżnimy samo istnienie takiej funkcji, jej niejednoznaczność oraz możliwość zapisania jej za pomocą znanych funkcji elementarnych.

\section{Definicja i jednoznaczność z dokładnością do stałej}
\begin{definicja}
\label{ma1:def:primitive}
Niech $U\subseteq\R$ będzie niepustym zbiorem otwartym, a $f:U\to\R$. Funkcję różniczkowalną $F:U\to\R$ nazywamy funkcją pierwotną funkcji $f$ na $U$, jeśli
\[
 F'(x)=f(x)\qquad(x\in U).
\]
\end{definicja}
W pierwszej kolejności będziemy rozpatrywać otwarte przedziały. Dla przedziału domkniętego można rozważać funkcję pierwotną ciągłą na całym przedziale, różniczkowalną w jego wnętrzu, z pochodnymi jednostronnymi na końcach, jeśli są one wymagane. W tej definicji i w następnych twierdzeniach nie używamy jednak pochodnych na brzegu zbioru otwartego.

\begin{twierdzenie}
\label{ma1:thm:primitive-uniqueness}
Niech $I$ będzie niepustym otwartym przedziałem i niech $F$ będzie funkcją pierwotną $f$ na $I$. Wszystkie funkcje pierwotne $f$ na $I$ mają postać
\[
 F+C,\qquad C\in\R.
\]
\end{twierdzenie}
\begin{proof}
Jeżeli $G$ jest drugą funkcją pierwotną $f$ na $I$, to z liniowości pochodnej (twierdzenie \ref{ma1:thm:derivative-algebra}) mamy $(G-F)'=G'-F'=0$. Dla dowolnych $u<v$ z $I$ odcinek $[u,v]$ zawiera się w $I$. Funkcja $G-F$ jest na nim ciągła i różniczkowalna w jego wnętrzu. Z twierdzenia Lagrange'a \ref{ma1:thm:lagrange} istnieje $c\in(u,v)$ takie, że
\[
 (G-F)(v)-(G-F)(u)=(v-u)(G-F)'(c)=0.
\]
Wobec dowolności $u,v$ funkcja $G-F$ jest stała na $I$. Odwrotnie, jeżeli $C\in\R$, to $(F+C)'=F'+0=f$, więc każda funkcja tej postaci jest pierwotna.
\end{proof}
Jest to także bezpośrednie zastosowanie wniosku \ref{ma1:cor:equal-derivatives}.

\begin{wniosek}[Warunek początkowy]
\label{ma1:cor:primitive-normalization}
Jeżeli $f$ ma funkcję pierwotną na przedziale $I$, to dla dowolnych $x_0\in I$ oraz $y_0\in\R$ istnieje dokładnie jedna funkcja pierwotna $G$ spełniająca $G(x_0)=y_0$.
\end{wniosek}
\begin{proof}
Jeżeli $F'=f$, to $G(x)=F(x)-F(x_0)+y_0$ spełnia oba warunki. Jeśli $G_1,G_2$ spełniają oba warunki, to $G_1-G_2$ jest stała z twierdzenia \ref{ma1:thm:primitive-uniqueness}; jej wartość w $x_0$ wynosi zero, więc $G_1=G_2$.
\end{proof}

\begin{center}
\begin{tikzpicture}[x=1.4cm,y=.85cm,>=Stealth]
 \draw[->] (-1.8,0)--(2.4,0) node[right] {$x$};
 \draw[->] (0,-1.5)--(0,5.5) node[above] {$y$};
 \draw[blue!65!black,thick,domain=-1.6:1.6,samples=90] plot (\x,{\x*\x});
 \draw[green!45!black,thick,domain=-1.6:1.6,samples=90] plot (\x,{\x*\x+2});
 \draw[orange!85!black,thick,domain=-1.6:1.6,samples=90] plot (\x,{\x*\x-1});
 \draw[gray,dashed] (1,-1.3)--(1,4.6);
 \fill[blue!65!black] (1,1) circle (1.7pt);
 \fill[green!45!black] (1,3) circle (1.7pt);
 \fill[orange!85!black] (1,0) circle (1.7pt);
 \draw[blue!65!black] (.65,.3)--(1.35,1.7);
 \draw[green!45!black] (.65,2.3)--(1.35,3.7);
 \draw[orange!85!black] (.65,-.7)--(1.35,.7);
 \node[right,blue!65!black] at (1.6,2.56) {$F(x)=x^2$};
 \node[right,green!45!black] at (1.6,4.56) {$F(x)+2$};
 \node[right,orange!85!black] at (1.6,1.56) {$F(x)-1$};
 \node[below right] at (1,-1.2) {$x_0=1$};
\end{tikzpicture}

\small Rysunek: trzy funkcje pierwotne funkcji $f(x)=2x$. W tym samym punkcie $x_0=1$ mają styczne o identycznym nachyleniu, lecz różne wartości. Warunek początkowy wybiera dokładnie jeden wykres z rodziny przesunięć pionowych.
\end{center}

\section{Dziedzina funkcji pierwotnej}
\begin{stwierdzenie}
\label{ma1:prop:primitive-components}
Niech $U$ będzie zbiorem otwartym w $\R$ i niech $F,G$ będą funkcjami pierwotnymi tej samej funkcji na $U$. Wtedy $G-F$ jest stała na każdym maksymalnym otwartym przedziale zawartym w $U$. Stałe na różnych takich przedziałach można wybierać niezależnie.
\end{stwierdzenie}
\begin{proof}
Dla $x\in U$ weźmy sumę wszystkich otwartych przedziałów zawartych w $U$ i zawierających $x$. Jest ona otwarta, zawarta w $U$ i jest przedziałem: jeśli $u<v$ należą do tej sumy, to należą do dwóch przedziałów zawierających wspólny punkt $x$. Suma tych dwóch przedziałów jest przedziałem, więc zawiera $[u,v]$. Otrzymany przedział jest maksymalny z konstrukcji. Dwa maksymalne przedziały nie mogą się przecinać, bo ich suma byłaby większym otwartym przedziałem zawartym w $U$. Każdy punkt $U$ należy więc do dokładnie jednego z nich.

Na każdym takim przedziale stosujemy twierdzenie \ref{ma1:thm:primitive-uniqueness}, otrzymując stałą różnicę. Odwrotnie, jeżeli na każdym przedziale $J$ wybierzemy liczbę $C_J$ i położymy $G(x)=F(x)+C_J$ dla $x\in J$, to każdy $x\in U$ ma otoczenie zawarte w swoim $J$. Na tym otoczeniu $G'=F'=f$, zatem $G$ jest funkcją pierwotną na $U$. Wybory $C_J$ są niezależne.
\end{proof}

\begin{przyklad}
\label{ma1:ex:primitive-log-domain}
Dla $f(x)=1/x$ na $\R\setminus\{0\}$ funkcja $F(x)=\log|x|$ jest pierwotna. Na $(0,\infty)$ jej pochodna to $1/x$ z twierdzenia \ref{ma1:thm:exp-log-derivatives}. Na $(-\infty,0)$ mamy $F(x)=\log(-x)$ i z reguły łańcuchowej \ref{ma1:thm:chain-rule} otrzymujemy $F'(x)=(-1)/(-x)=1/x$. Wszystkie funkcje pierwotne na całej dziedzinie mają więc postać
\[
 G(x)=
 \begin{cases}
 \log(-x)+C_-,&x<0,\\
 \log x+C_+,&x>0,
 \end{cases}
 \qquad C_-,C_+\in\R.
\]
Nie ma powodu, aby $C_-=C_+$. Na przykład wybór $C_-=0$, $C_+=1$ daje funkcję pierwotną, której nie można zapisać jako $\log|x|+C$ z jedną globalną stałą.
\end{przyklad}

\section{Całka nieoznaczona i liniowość}
\begin{definicja}
\label{ma1:def:indefinite-integral}
Jeżeli $f:I\to\R$ ma funkcję pierwotną $F$ na niepustym otwartym przedziale $I$, to zapis
\[
 \int f(x)\,\dd x=F(x)+C,\qquad C\in\R,
\]
oznacza rodzinę wszystkich jej funkcji pierwotnych na $I$. Nazywamy ją całką nieoznaczoną funkcji $f$ na $I$.
\end{definicja}
Symbol $\dd x$ wskazuje zmienną całkowania. Wiąże się z różniczką z definicji \ref{ma1:def:differential}: jeśli $F'=f$, to $\dd F_x=f(x)\dd x_x$ jako równość odwzorowań liniowych w punkcie $x$. Zapis całki nieoznaczonej nie oznacza jeszcze liczby, pola powierzchni ani granicy sum. W następnym rozdziale zdefiniujemy całkę oznaczoną Riemanna i wykażemy jej związek z funkcjami pierwotnymi.

\begin{twierdzenie}
\label{ma1:thm:primitive-linearity}
Jeżeli $F'=f$ i $G'=g$ na przedziale $I$, to dla dowolnych $\alpha,\beta\in\R$ wszystkie funkcje pierwotne $\alpha f+\beta g$ na $I$ mają postać
\[
 \alpha F+\beta G+C,\qquad C\in\R.
\]
\end{twierdzenie}
\begin{proof}
Z liniowości pochodnej mamy $(\alpha F+\beta G)'=\alpha f+\beta g$. Funkcja $\alpha F+\beta G$ jest więc jedną funkcją pierwotną, a twierdzenie \ref{ma1:thm:primitive-uniqueness} opisuje wszystkie pozostałe.
\end{proof}
W obliczeniach często zapisujemy liniowość symbolicznie jako $\int(\alpha f+\beta g)\,\dd x=\alpha\int f\,\dd x+\beta\int g\,\dd x$. Należy wówczas rozumieć prawą stronę jako wybór funkcji pierwotnych i dodanie dowolnej stałej końcowej. Jeżeli $\alpha=\beta=0$, całka lewej strony oznacza wszystkie funkcje stałe, a dosłowne mnożenie dwóch rodzin przez zero i ich dodawanie dałoby jedynie funkcję zerową. Wersja twierdzenia \ref{ma1:thm:primitive-linearity} jest wolna od tej niejednoznaczności.

\section{Podstawowe funkcje pierwotne}
\begin{stwierdzenie}
\label{ma1:prop:primitive-table}
Poniższe wzory podają po jednej funkcji pierwotnej. Aby otrzymać wszystkie funkcje pierwotne na danym otwartym przedziale, należy dodać dowolną stałą.
\[
\begin{array}{c|c|l}
 f(x)&F(x)&\text{Dziedzina wzoru}\\ \hline
 1 & x & \R\\
 x^r,\ r\in\R\setminus\{-1\} & x^{r+1}/(r+1)& x>0\\
 1/x & \log|x| & x\ne0\\
 e^x & e^x & \R\\
 a^x,\ a>0,\ a\ne1 & a^x/\log a & \R\\
 \cos x & \sin x & \R\\
 \sin x & -\cos x & \R\\
 1/\cos^2 x & \tan x & \cos x\ne0\\
 1/\sin^2 x & -\cot x & \sin x\ne0\\
 1/(1+x^2) & \arctan x & \R\\
 1/\sqrt{1-x^2} & \arcsin x & |x|<1
\end{array}
\]
Dla całkowitego $r\ge0$ wzór potęgowy zachodzi na całej $\R$, a dla całkowitego $r<-1$ na $\R\setminus\{0\}$. W przypadku $r=0$ funkcję $x^0$ rozumiemy jako wielomian stały $1$.
\end{stwierdzenie}
\begin{proof}
Dla $x>0$ różniczkujemy $x^{r+1}=\exp((r+1)\log x)$, otrzymując
\[
 \left(\frac{x^{r+1}}{r+1}\right)'
 =\frac{\exp((r+1)\log x)}{r+1}\frac{r+1}{x}
 =x^r.
\]
Dla całkowitych wykładników na wskazanych większych dziedzinach stosujemy wzór \ref{ma1:cor:integer-power-derivative}. Wiersz $1/x$ uzasadniliśmy w przykładzie \ref{ma1:ex:primitive-log-domain}. Dla $a^x=\exp(x\log a)$ mamy $(a^x)'=(\log a)a^x$ z reguły łańcuchowej \ref{ma1:thm:chain-rule}; $\log a\ne0$ dla $a\ne1$. Pozostałe wiersze są bezpośrednio wzorami z twierdzeń \ref{ma1:thm:exp-log-derivatives}, \ref{ma1:thm:trig-derivatives} i \ref{ma1:thm:inverse-trig-derivatives}, odczytanymi jako równości $F'=f$. Jeżeli dziedzina wzoru jest rozłączną sumą przedziałów, opis wszystkich funkcji pierwotnych daje stwierdzenie \ref{ma1:prop:primitive-components}.
\end{proof}

Dodatkowe wzory wynikające z twierdzenia \ref{ma1:thm:hyperbolic-derivatives}:
\[
\begin{array}{c|c|l}
 f(x)&F(x)&\text{Dziedzina wzoru}\\ \hline
 \cosh x & \sinh x & \R\\
 \sinh x & \cosh x & \R\\
 1/\cosh^2x & \tanh x & \R\\
 1/\sinh^2x & -\coth x & x\ne0\\
 1/\sqrt{1+x^2} & \operatorname{arsinh}x & \R\\
 1/\sqrt{x^2-1} & \operatorname{arcosh}x & x>1\\
 1/(1-x^2) & \operatorname{artanh}x & |x|<1\\
 1/(1-x^2) & \operatorname{arcoth}x & |x|>1
\end{array}
\]
Na $x<-1$ funkcją pierwotną $1/\sqrt{x^2-1}$ jest $-\operatorname{arcosh}(-x)$: reguła łańcuchowa daje pochodną $-(-1)/\sqrt{(-x)^2-1}=1/\sqrt{x^2-1}$.

\begin{przyklad}
\label{ma1:ex:primitive-polynomial}
Na $\R$ szukamy funkcji pierwotnej $f(x)=3x^2-4x+5$ spełniającej $G(1)=2$. Z tabeli i liniowości otrzymujemy $F(x)=x^3-2x^2+5x$. Ponieważ $F(1)=4$, warunek początkowy daje $G(x)=F(x)-4+2=x^3-2x^2+5x-2$. Różniczkowanie potwierdza $G'=f$, a podstawienie $x=1$ daje $G(1)=2$. Jednoznaczność wynika z wniosku \ref{ma1:cor:primitive-normalization}.
\end{przyklad}

\begin{przyklad}
\label{ma1:ex:primitive-affine-argument}
Dla $a,b\in\R$, $a\ne0$, jeśli $F'=f$ na przedziale $I$, to na przedziale $J=\{x:ax+b\in I\}$ funkcja $H(x)=F(ax+b)/a$ jest funkcją pierwotną $x\mapsto f(ax+b)$. Istotnie,
\[
 H'(x)=\frac1a F'(ax+b)a=f(ax+b)
\]
z reguły łańcuchowej. W szczególności na $\R$ mamy
\[
 \int \cos(3x-1)\,\dd x=\frac{\sin(3x-1)}3+C.
\]
Nie trzeba osobno zgadywać współczynnika $1/3$: wynika on z konieczności skompensowania pochodnej argumentu. Ogólną metodę podstawienia i całkowanie przez części omówimy w rozdziale poświęconym metodom całkowania.
\end{przyklad}

\section{Istnienie funkcji pierwotnej}
\begin{stwierdzenie}
\label{ma1:prop:primitive-necessary}
Jeżeli $f$ ma funkcję pierwotną na otwartym przedziale $I$, to $f$ ma własność wartości pośredniej na $I$. Ponadto, jeśli w punkcie $a\in I$ istnieje jednostronna granica $f$ w rozszerzonej prostej rzeczywistej, to jest skończona i równa $f(a)$.
\end{stwierdzenie}
\begin{proof}
Jeżeli $F'=f$, to pierwsza teza wynika z twierdzenia Darboux \ref{ma1:thm:derivative-darboux} zastosowanego do $F$. Drugą tezę otrzymujemy z wniosku \ref{ma1:cor:derivative-no-jumps}.
\end{proof}
\begin{przyklad}
\label{ma1:ex:primitive-nonexistence}
Funkcja $f(x)=0$ dla $x<0$ oraz $f(x)=1$ dla $x\ge0$ nie ma funkcji pierwotnej na $\R$. Gdyby $F'=f$, to z $F'(-1)=0$ i $F'(1)=1$ wynikałoby istnienie $c\in(-1,1)$ z $F'(c)=1/2$, wbrew temu, że $f$ przyjmuje jedynie wartości $0,1$. Na każdej z dwóch półprostych otwartych funkcja ma jednak pierwotną, odpowiednio stałą lub postaci $x+C$.
\end{przyklad}

\begin{przyklad}
\label{ma1:ex:primitive-discontinuous}
Ciągłość nie jest warunkiem koniecznym istnienia funkcji pierwotnej. W przykładzie \ref{ma1:ex:discontinuous-derivative} wykazaliśmy, że funkcja
\[
 F(x)=
 \begin{cases}x^2\sin(1/x),&x\ne0,\\0,&x=0\end{cases}
\]
jest różniczkowalna na $\R$, a jej pochodna
\[
 f(x)=
 \begin{cases}2x\sin(1/x)-\cos(1/x),&x\ne0,\\0,&x=0\end{cases}
\]
nie jest ciągła w zerze. Mimo to $F$ jest funkcją pierwotną $f$ na całej $\R$.
\end{przyklad}

W kolejnym rozdziale udowodnimy podstawowe twierdzenie rachunku całkowego: dla funkcji ciągłej $f$ na otwartym przedziale $I$ oraz ustalonego $x_0\in I$ funkcja
\[
 F(x)=\int_{x_0}^{x} f(t)\,\dd t
\]
jest pierwotna $f$ i spełnia $F(x_0)=0$. Jest to zapowiedź wyniku, nie zastosowanie niezdefiniowanej jeszcze całki oznaczonej. Dopiero po skonstruowaniu całki Riemanna uzyskamy dowód istnienia funkcji pierwotnej dla każdej funkcji ciągłej. Z samego istnienia nie wynika, że funkcję pierwotną można wyrazić za pomocą skończonej liczby działań i znanych funkcji elementarnych; zagadnienie takiej reprezentacji jest odrębne.
